\section{Part 2: entrenamiento de modelos «Thinking»}

\begin{figure}[H]
\centering
\includegraphics[width=0.9\textwidth]{slides-images/slide-30.jpg}
\caption{Part 2: principales artículos sobre el entrenamiento de modelos «Thinking»\protect\footnotemark}
\end{figure}
\footnotetext{Slides, página 30.}

\subsection{DeepSeek-R1: incentivar la capacidad de razonamiento mediante RL}

DeepSeek-R1 es el trabajo emblemático del paradigma del «thinking model». Su hallazgo central es: \textbf{solo con entrenamiento por RL (sin cadenas de razonamiento anotadas por personas), el modelo desarrolla de forma espontánea comportamientos de razonamiento complejos}.

\begin{importantbox}{Flujo de entrenamiento de DeepSeek-R1}
\begin{enumerate}
    \item Se aplica GRPO (Group Relative Policy Optimization) directamente sobre el modelo base para el entrenamiento por RL
    \item La recompensa se basa únicamente en si la respuesta final es correcta (outcome reward), más una recompensa de formato
    \item El modelo aprendió de forma espontánea metaestrategias como la verificación (verification), la vuelta atrás (backtracking), el establecimiento de submetas (subgoal setting) y el razonamiento hacia atrás (backward chaining)
\end{enumerate}
\end{importantbox}

\begin{knowledgebox}{GRPO (Group Relative Policy Optimization)}
GRPO es una variante simplificada de PPO. Para cada problema se muestrea un grupo de soluciones y se usa la recompensa relativa dentro del grupo (restando la media del grupo y normalizando) como estimación de la ventaja, lo que evita entrenar una red critic/value adicional. Esto resulta más eficiente en el entrenamiento de LLM a gran escala.
\end{knowledgebox}

\subsection{Entender qué es el «thinking» a partir de trayectorias concretas}

La expresión «Thinking model» se presta fácilmente a descripciones misteriosas, pero los ejemplos que muestra el ponente son en realidad muy concretos: dentro de una sola respuesta, el modelo ejecuta más operaciones de nivel macro, por ejemplo, probar una dirección, descubrir que no funciona, verificar de forma explícita, retroceder y volver a planificar. Es decir, el RL no se limita a alargar la salida, sino que mejora la capacidad del modelo para organizar estas acciones de nivel macro dentro de un presupuesto fijo.

\begin{figure}[H]
\centering
\includegraphics[width=0.9\textwidth]{slides-images/slide-31.jpg}
\caption{Esquema de una sola trayectoria de un Thinking model: la verificación, la vuelta atrás y la replanificación se alternan\protect\footnotemark}
\end{figure}
\footnotetext{Slides, página 31.}

\subsection{Lo que realmente cambia es el espacio de acciones, no el libro de texto de RL}

En las slides, el ponente presenta de forma expresa un juicio muy importante: \textbf{el objetivo de entrenamiento no ha cambiado en lo esencial; lo que realmente cambió son las «macro actions» que el modelo ya puede ejecutar.} DeepSeek-R1 sigue haciendo policy gradient, y Kimi K1.5 también sigue dependiendo de actualizaciones de la política impulsadas por la advantage; pero el modelo base ya puede expresar, con un token budget más largo, comportamientos de alto nivel como la verificación, la vuelta atrás y la descomposición en submetas, por lo que la misma receta de RL puede producir efectos más fuertes en el nuevo espacio de acciones.

\begin{figure}[H]
\centering
\includegraphics[width=0.9\textwidth]{slides-images/slide-32.jpg}
\caption{Cambios en la era de los Thinking models: la forma de la recompensa permanece casi igual, mientras que el espacio de acciones y el presupuesto se amplían de manera notable\protect\footnotemark}
\end{figure}
\footnotetext{Slides, página 32.}

\begin{knowledgebox}{Por qué «mejorar el espacio de acciones» es más decisivo que «mejorar la función objetivo»}
Muchas discusiones atribuyen el avance de los thinking models a algún optimizador nuevo o a una nueva pérdida, pero el juicio de esta clase es más mesurado: basta con que el modelo base ya tenga cierto germen de metacognición para que ampliar el presupuesto de contexto, permitir rollouts más largos y luego recompensar estos comportamientos con RL amplifique de forma notable el desempeño en razonamiento. Dicho de otro modo, \textbf{las capacidades nuevas suelen surgir de amplificar y recombinar capacidades existentes, no de inventar de la nada un objetivo de entrenamiento completamente nuevo.}
\end{knowledgebox}

\subsection{Comportamientos de razonamiento emergentes}

\begin{figure}[H]
\centering
\includegraphics[width=0.9\textwidth]{slides-images/slide-35.jpg}
\caption{«Action» Space: comparación de las metaestrategias emergentes\protect\footnotemark}
\end{figure}
\footnotetext{Slides, página 35.}

El estudio de Gandhi et al. (2025) muestra que el que surjan comportamientos cognitivos durante el entrenamiento por RL depende de si el modelo base ya posee las «semillas» de esos comportamientos. En los modelos de la serie DeepSeek, durante el entrenamiento por RL, la frecuencia de los comportamientos de verificación y de vuelta atrás primero aumenta y después disminuye; en cambio, los modelos Llama de Meta no mostraron en absoluto estos comportamientos, y el entrenamiento por RL tampoco logró mejorar su capacidad de razonamiento.

\begin{warningbox}{El RL no lo resuelve todo}
El RL no puede hacer que un modelo aprenda comportamientos que «no sabe hacer en absoluto». Si en la etapa de preentrenamiento el modelo base no estuvo expuesto a datos de razonamiento en cadena, es muy difícil que solo con RL surjan desde cero comportamientos de razonamiento complejos. La capacidad del modelo base es la condición previa para que el RL tenga efecto.
\end{warningbox}

\subsection{Objetivo de entrenamiento con presupuesto fijo}

A continuación, el ponente reformula el problema como un objetivo muy de ingeniería: dado que cada pregunta de prueba cuenta con un compute budget total fijo, se entrena una política para que la sampled response maximice la tasa de éxito dentro de ese presupuesto. Este planteamiento tiene dos implicaciones prácticas:

\begin{enumerate}
    \item Razonar no consiste en que «más largo es mejor», sino en decidir, dentro de un \textbf{presupuesto limitado}, en qué gastar el cómputo.
    \item RL, SFT y RFT pueden verse, en esencia, como distintos métodos de solución de este problema de adaptación; sus diferencias residen principalmente en si aprovechan la recompensa y en qué tan a fondo la aprovechan.
\end{enumerate}

\begin{figure}[H]
\centering
\includegraphics[width=0.9\textwidth]{slides-images/slide-33.jpg}
\caption{Objetivo del entrenamiento de Thinking models: maximizar la tasa de éxito dentro de un presupuesto de prueba fijo\protect\footnotemark}
\end{figure}
\footnotetext{Slides, página 33.}

\begin{figure}[H]
\centering
\includegraphics[width=0.9\textwidth]{slides-images/slide-34.jpg}
\caption{Bajo la misma restricción de presupuesto, tanto RL como SFT/RFT pueden verse como implementaciones distintas del «problema de adaptación»\protect\footnotemark}
\end{figure}
\footnotetext{Slides, página 34.}

\subsection{Kimi K1.5: escalar el entrenamiento por RL}

La contribución central de Kimi K1.5 consiste en cómo realizar de manera eficaz el entrenamiento por RL a gran escala:
\begin{itemize}
    \item Una ventana de contexto más larga (128K tokens)
    \item Mejores estrategias de muestreo y un mejor diseño de la recompensa
    \item Los experimentos muestran que el desempeño del entrenamiento por RL sigue mejorando a medida que aumenta el cómputo
\end{itemize}

\subsection{Amplificación de la longitud y el nuevo paradigma de test-time scaling}

Otra consecuencia directa del entrenamiento por RL es que el modelo está dispuesto a dedicar más presupuesto de tokens a lo que realmente aporta valor. Citando el informe de Kimi, el ponente señala que, tras el entrenamiento por RL, la longitud de las respuestas del modelo crece de forma notable; pero no se trata de una verbosidad sin sentido, sino que va acompañada de más verificaciones, correcciones y exploración de ramas. Así, «dar un poco más de cómputo en el momento de la inferencia» empieza a convertirse en una vía de mejora del desempeño que puede situarse a la par de «seguir aumentando los parámetros del modelo».

\begin{figure}[H]
\centering
\includegraphics[width=0.9\textwidth]{slides-images/slide-36.jpg}
\caption{El entrenamiento por RL amplifica la longitud de las respuestas y convierte esa longitud adicional en comportamientos de razonamiento más ricos\protect\footnotemark}
\end{figure}
\footnotetext{Slides, página 36.}

\begin{figure}[H]
\centering
\includegraphics[width=0.9\textwidth]{slides-images/slide-37.jpg}
\caption{Nuevo planteamiento del problema: además de escalar el modelo, también se puede escalar el test-time compute\protect\footnotemark}
\end{figure}
\footnotetext{Slides, página 37.}

\begin{figure}[H]
\centering
\includegraphics[width=0.9\textwidth]{slides-images/slide-39.jpg}
\caption{Un pensamiento más largo y una autocorrección más fuerte conforman la nueva vía de test-time scaling\protect\footnotemark}
\end{figure}
\footnotetext{Slides, página 39.}

\subsection{Resultados de los modelos Thinking}

\begin{figure}[H]
\centering
\includegraphics[width=0.9\textwidth]{slides-images/slide-40.jpg}
\caption{Desempeño de los modelos Thinking en distintos benchmarks\protect\footnotemark}
\end{figure}
\footnotetext{Slides, página 40.}

Los modelos Thinking lograron mejoras notables en varios bancos de prueba, como la programación competitiva de Codeforces, Humanity's Last Exam y AIME. DeepSeek-R1 alcanza un Pass@1 de 79.8\% en AIME 2024, comparable con el 79.2\% de OpenAI o1-1217.

\subsection{Test-Time Compute y Meta-RL}

Por último, el ponente menciona que optimizar el test-time compute (la asignación de cómputo en el momento de la inferencia) es, en esencia, un problema de Meta-RL: el modelo debe aprender cómo «gastar» el presupuesto de cómputo durante la inferencia.

\subsection{Problemas aún sin resolver}

Aunque los resultados en las leaderboards ya son muy llamativos, la última página de esta clase es, por el contrario, la que tiene el tono de advertencia más fuerte: los thinking models actuales todavía carecen de una formulation unificada y clara. Aún no hemos respondido por completo tres cuestiones: primero, si existe un diseño de dense reward más robusto que el outcome reward; segundo, en qué condiciones la ganancia del RL supera de manera estable a la del SFT; tercero, cómo enunciar con rigor el «reforzamiento del razonamiento» como un problema de adaptación, y no como una acumulación de técnicas empíricas.

\begin{figure}[H]
\centering
\includegraphics[width=0.9\textwidth]{slides-images/slide-41.jpg}
\caption{Preguntas abiertas planteadas al cierre del curso: la formulation, las dense rewards y los límites de la ventaja del RL sobre el SFT\protect\footnotemark}
\end{figure}
\footnotetext{Slides, página 41.}

\begin{warningbox}{El verdadero reto no está en replicar un DeepSeek-R1}
Si uno se fija solo en algún modelo público o en alguna curva de benchmark, es fácil malinterpretar el problema como «encontrar la misma recipe». Pero la manera en que Aviral Kumar cierra la clase se parece más a una lista de preguntas de investigación: todavía necesitamos esclarecer la forma de la señal de recompensa, el modo en que se acoplan la búsqueda y la verificación, y por qué los distintos modelos base responden al RL de maneras tan diferentes.
\end{warningbox}

\subsection{Resumen de la sección}

Los thinking models como DeepSeek-R1 demostraron que: (1) el RL puede incentivar que los LLM desarrollen por sí mismos capacidades de razonamiento; (2) la condición previa decisiva es que el modelo base posea las «semillas» de los comportamientos de razonamiento; (3) los métodos simplificados de optimización de políticas, como GRPO, son a la vez eficientes y eficaces en el entrenamiento a gran escala.
